% !TEX program = xelatex
\documentclass[12pt,a4paper]{article}

\usepackage[a4paper,margin=1in]{geometry}
\usepackage{fontspec}
\usepackage{microtype}
\usepackage{setspace}
\usepackage{booktabs}
\usepackage{tabularx}
\usepackage{array}
\usepackage[hidelinks]{hyperref}

\setmainfont{Times New Roman}
\onehalfspacing
\setlength{\parindent}{1.5em}
\setlength{\parskip}{0.35em}
\setcounter{secnumdepth}{2}
\setcounter{tocdepth}{2}

\newcolumntype{Y}{>{\raggedright\arraybackslash}X}

\title{\textbf{Literature Review}}
\author{}
\date{}

\begin{document}

\maketitle

\section{Introduction}

\subsection{Background}

As higher education institutions face an increasingly diverse student body and the demand for more accessible and personalised education, chatbots emerge as an innovative solution to address a myriad of challenges \cite{martinez2024}.

When students move from high school to an unfamiliar university environment, they will naturally have many related questions to ask. Although university students are increasingly familiar with using the internet to obtain information, the information they need is often scattered across different websites, documents, and communication platforms. As a result, finding accurate and timely answers might be challenging.

In the paper by Yuwei Qian, it is pointed out that information-seeking behavior varies across individuals in different environments. This is especially true for freshmen, who may find it difficult to begin a university life that is very different from their high school experience \cite{qian2019}. This highlights the importance of providing structured and easily accessible information systems for freshmen.

Moreover, due to different living habits in different countries, students use a wide variety of online communication software. This also makes it difficult for universities to fully deliver daily-life information beyond official announcements.

Therefore, there is a need for a unified chatbot that can provide structured and reliable access to university information for freshmen. And in order to build the chatbot, a systematic review on previous chatbot applications is needed.

\subsection{Purpose and Scope}

The purpose of this literature review is to examine chatbot applications developed for university students, especially freshmen. It explores the previous applications of chatbots in the domain of education, the approaches adopted in these systems, and the evaluations metrics used for different types of chatbots, in order to better guide the development of the chatbot.

This review focuses on chatbot systems designed for higher education contexts, particularly those that support university students and freshmen through university enquiry services, student support, and FAQ-based assistance. It mainly examines traditional chatbot approaches, including rule-based, machine-learning-based, retrieval-based, and hybrid methods, while modern large language model (LLM)-based generative chatbots are discussed only briefly rather than in detail. The review also covers the ways chatbot systems have been evaluated in earlier studies. Among the measures most often reported are classification accuracy, retrieval performance, response quality, and user satisfaction.

A comparison of these studies shows the different approaches used in university chatbot development and evaluation. It also draws attention to a few issues that remain open for further investigation.

\subsection{Research Questions}

\begin{description}
    \item[\textbf{RQ1.}] What existing chatbot applications have been developed to support university students?
    \item[\textbf{RQ2.}] What techniques have been used in chatbot systems?
    \item[\textbf{RQ3.}] What evaluation methods are used to assess diffenrent types of chatbots?
\end{description}

\subsection{Search Strategy and Selection Criteria}

The sources used in this study were mainly obtained from academic databases, including Google Scholar and IEEE Xplore. To develop a broader understanding of university chatbot research, journal articles, conference papers, and selected academic books were reviewed. Most of the sources reviewed were published between 2008 and 2026. This period covers both early studies on university chatbots and more recent research influenced by developments in artificial intelligence.

Relevant materials were identified through keyword searches and reference checking. Keywords used in the search included “university counseling chatbot,” “higher education chatbot,” “student support chatbot,” “FAQ chatbot,” “rule-based chatbot,” “retrieval-based chatbot,” “hybrid chatbot,” and “information retrieval in chatbots.” After an initial group of papers had been selected, the reference lists of several key studies were examined to identify additional publications related to the topic. This process helped expand the range of materials included in the review.

As the project is focused on higher education, particular attention was given to studies on student support services, campus information systems, and FAQ-based chatbot applications. Research related to rule matching, machine learning classification, and information retrieval was also included because these techniques are relevant to the chatbot developed in this project.

Studies related mainly to e-commerce, customer service, and other non-educational settings were used only for background reference. Papers that discussed chatbot concepts without presenting system development, implementation, or evaluation were also given less emphasis. The selected references provide useful information on current university chatbot applications, commonly used technologies, and evaluation methods. These studies also support the discussion of the research questions addressed in this project.
\section{Existing Chatbot Applications for University Students and Freshmen}

\subsection{University Enquiry Chatbots}

University enquiry chatbots are widely used in higher education to answer student questions and provide access to campus information. Common uses include administrative enquiries, university announcements, and information about student services.

Susanna et al.\ introduced a College Enquiry Chatbot that allows students to obtain information from links, text messages, and PDF files through one platform \cite{susanna2020}. Salve et al.\ developed a chatbot for a similar purpose. Their system also allows students to provide feedback when an answer is incorrect, making it possible for administrators to review and update the stored information \cite{salve}.

The two studies share a similar goal but use different approaches. Susanna et al.\ concentrate on bringing information from different sources into one system. Salve et al.\ focus more on improving answer quality through user feedback. The studies suggest that easy access to information alone is not enough. The information also needs to remain accurate and updated so that students can continue to receive useful responses.

\subsection{Student Support and FAQ Chatbots}

Student support and FAQ chatbots generally cover a wider range of functions than enquiry chatbots. Instead of only providing information, they are often designed to assist communication, learning activities, and student services.

Dhandayuthapani introduced a student support chatbot framework that supports multiple languages and different communication methods, including text, voice, and graphical interfaces \cite{dhandayuthapani2022}. The study suggests that such systems may improve service delivery while helping to reduce staff workload. Nourbakhsh and Yaser also discussed the role of chatbots in education, arguing that they can support learning, improve communication between students and teachers, and make educational resources more accessible \cite{nourbakhsh2017}.

The studies reviewed show that higher education chatbots are no longer limited to answering questions or providing campus information. Increasingly, they are being developed to support students in different areas of university life. This development reflects a growing interest in using chatbots not only as information tools but also as part of the broader student support system.

Although many studies describe the benefits of these systems, much of the existing work focuses on proposed models and expected benefits rather than real-world implementation. More research is therefore needed to examine how these chatbots operate in actual university settings and how students respond to them during regular use.

\subsection{Limitations of Existing University Chatbot Applications}

Previous studies have identified several limitations of university chatbots. According to Raval, chatbots may struggle to understand slang, informal language, and messages that contain strong emotions. When this happens, user requests may be interpreted incorrectly, or the chatbot may repeat the same response several times during a conversation \cite{raval2020}. These problems can reduce the accuracy of chatbot responses and may lead to a less positive user experience.

A further difficulty comes from the nature of university information. Different institutions follow different policies, procedures, and administrative practices. As a result, chatbot databases need to be updated regularly and supported by a large amount of information. Keeping the content correct and up to date can be difficult, particularly when institutional information changes over time.

Many chatbot applications also rely on predefined knowledge bases and prepared responses. While this approach is suitable for handling common enquiries, it is less effective when users ask unclear questions or requests that fall outside expected patterns. Much of the existing research focuses on system features and technical design. Fewer studies have examined user satisfaction over extended periods or evaluated system performance in everyday use.

Based on the studies reviewed, no single chatbot method appears capable of meeting all student needs. Rule-based approaches, machine learning techniques, and information retrieval methods each perform well in certain areas but also have limitations. For this reason, combining these techniques within a hybrid chatbot may offer a more balanced approach to student support.

\section{Technical Approaches Used in University Chatbots}

In the development of university FAQ chatbots, it is essential to notice the technical architecture because it directly determines the overall performance of systems. To address the diverse needs of students, modern conversational agents have evolved from single method frameworks into multi-modular systems. This section evaluates various methods and deeply analyzes the core text classification algorithms supporting intent routing.

\subsection{Rule-Based Approaches}

Rule-based chatbots represent the earliest chatbot paradigm in educational systems. They provide specific responses relied on predefined rules and patterns for matching user inputs. Among the reviewed studies, we find Artificial Intelligence Markup Language (AIML) is commonly used to create rule-based chatbots. Khin and Soe \cite{khin2020}, for example, developed a university chatbot using AIML to provide predefined responses to frequently asked questions. Since they are easy to build, developers tend to use them to address well-defined queries, such as assignments distribution or course schedules \cite{nguyen2021,alghadhban2020,dinesh2021}, which can greatly reduce the effort and cost to build other chatbots. However, Lopez and Qamber \cite{lopez2022} observed that chatbots in higher education may struggle when users employ unexpected expressions or submit complex and ambiguous questions. Therefore, the reviewed literature suggests that their reliance on fixed patterns and short for semantic understanding makes them appropriate for narrowly defined queries, which are generally insufficient for university chatbot.

\subsection{Machine-Learning-Based Approaches}

To overcome the rigidity of pattern matching, researchers have increasingly used machine-learning-based (ML-based) methods for intent and category classification. As the development of natural language processing and machine learning algorithms, the methods for extracting text feature have been thriving. In the text classification task, the term frequency-inverse document frequency (TF-IDF) algorithm is widely used for its simplicity, strong interpretability, and suitability for small and medium-sized corpora. Its simple structure and clearing feature meaning make it adapt to constrained dialogue system in FAQ chatbots \cite{ashraf2024,yang2026}. Based on the text feature extraction, we evaluate three core text classification algorithms:

\begin{itemize}
    \item \textbf{Naive Bayes:} Naive Bayes is a probabilistic machine learning algorithm that can be used in a wide variety of classification tasks. In early text classification, the multivariate Bernoulli model only recorded the presence or absence of a word, ignoring the term frequency. A study found that Multinomial Naive Bayes (MNB) model, which incorporates word frequency, performs better at larger vocabulary sizes \cite{mccallum1998}. However, an unseen word may produce a zero probability if it does not occur in the training data. To addressed this issue, Setyaningsih and Listiowarni \cite{setyaningsih2021} applied Laplace smoothing to ensure the robustness of inputs. And it is also proved to minimize misclassification errors. Consequently, MNB with smoothing provides a efficient and relatively robust solution for intent classification, especially when the available training dataset is limited.

    \item \textbf{Logistic Regression:} Among the papers we collected, Logistic Regression is another widely used method. Unlike MNB fitting a joint probability distribution, it directly models the posterior probability, which maps from features to class labels. It can model the relationship between TF-IDF features and intent categories. However, Acito pointed out that some rare or unseen words may be assigned extreme weights, which can cause overfitting, misleading models in real world scenario \cite{acito2024}. To overcome this problem, many developers apply L2 regularization to improve generalizability and mitigate overfitting.

    \item \textbf{Linear SVM:} SVM is also a common geometric classification method based on structural risk minimization. It searches the best compromise between the complexity of the model (the learning accuracy of a particular training sample) and learning ability (error-free ability to identify any samples) in order to expect the best generalization ability \cite{li2009}. Liu et al.\ compared SVM with several other text classification methods and found that SVM achieved stronger text classification performance in their experiment \cite{liu2010}. These findings predicate that Linear SVM is particularly suitable for high-dimensional and sparse textual features.
\end{itemize}

Overall, the reviewed studies demonstrate that MNB, Logistic Regression, and Linear SVM are all viable algorithms for intent classification. MNB can save computational sources and it also performs well with word-frequency features, Logistic Regression provides probabilistic outputs, and Linear SVM often excels in high-dimensional feature spaces. However, the cited studies were conducted on different datasets and tasks. Their reported numerical results therefore cannot be compared directly. For the proposed chatbot, all three algorithms should be evaluated using the same university-query dataset.

\subsection{Retrieval-Based Approaches}

A retrieval-based chatbot works by matching user inputs to the most relevant responses within a predefined dataset. The responses here are entered manually by the developer, or based on a knowledge base of pre-existing information. Traditional systems commonly use lexical retrieval methods such as TF-IDF with cosine similarity or BM25 for their efficiency and interpretability \cite{salton1988,robertson2009}. However, these methods may perform poorly when users express the same meaning with different vocabulary. In order to solve this limitation, researchers introduced semantic retrieval methods based on BERT, Sentence-BERT, and dense passage representations, which can better identify paraphrases and semantically related questions \cite{sakata2019,reimers2019,karpukhin2020}. Although these approaches improve retrieval flexibility, their effectiveness still depends heavily on the quality of the knowledge base. Therefore, retrieval-based methods are particularly suitable for university service chatbots because they provide reviewed and traceable answers to official enquiries.

\subsection{Hybrid Chatbot Architectures}

Among the papers we collected, quite a lot researchers choose Hybrid Chatbot Architecture to combine the rule-based systems, the ML models and retrieval-based or generative components. In a hybrid architecture, an intent classification model is first used to identify the type of user query and then direct it to the most appropriate module. Researchers such as Mikael et al.\ \cite{mikael2025} proposed a hybrid chatbot for educational administration. They used a MNB classifier for query routing. Questions related to registration or financial matters were handled by the rule-based component, while academic and campus-related questions were directed to a retrieval-based module. For more complex or open-ended queries, they are forwarded to the generative model. By combining these different approaches, the system can improve both response reliability and conversational flexibility.

The reviewed hybrid studies show that different modules can be assigned according to their respective strengths. However, hybrid architectures are more difficult to design and maintain. An incorrect intent prediction may direct a query to an unsuitable module even when the required information is available elsewhere in the system. As a result, it is important to evaluate both classifier performance and end-to-end chatbot performance.

\subsection{Comparison of Methods}

To explicitly justify the technical selection, Table summarizes the strengths, limitations, and suitability of the discussed chatbot approaches.

\begin{table}[htbp]
    \centering
    \caption{Comparison of chatbot approaches}
    \label{tab:method-comparison}
    \renewcommand{\arraystretch}{1.25}
    \small
    \begin{tabularx}{\textwidth}{>{\raggedright\arraybackslash}p{0.16\textwidth}YYY}
        \toprule
        \textbf{Approach} & \textbf{Strength} & \textbf{Limitation} & \textbf{Suitability for Our Project} \\
        \midrule
        Rule-based & Controllable and simple & Poor flexibility & Good for greetings/help \\
        ML-based & Handles varied wording & Needs labelled data & Good for category classification \\
        Retrieval-based & Uses reviewed answers & Depends on FAQ coverage & Good for official FAQ answers \\
        Generative & Flexible & Risk of hallucination & Not main method \\
        Hybrid & Balances control and flexibility & More complex & Most suitable \\
        \bottomrule
    \end{tabularx}
\end{table}

Generally, each chatbot approach has its own strengths and limitations. Since hybrid approach combines control, flexibility, and reliability, it is the most common architecture for university chatbots.

\section{Evaluation Methods for Chatbots}

Since different types of chatbots have their own characteristics, they cannot be evaluated using the same method. 
Therefore, a systematic review of evaluation metrics for different types of chatbots is needed.

\subsection{Evaluation Metrics for Machine Learning Category Classification}

With the emergence of different types of chatbots, machine learning is widely used to allocate tasks to chatbots with different strengths. And evaluating the classification becomes a need. Luckily, machine learning is quite a mature technique, and most of the classic evaluation metrics for machine learning can be directly applied here. In the 34 papers we collected, 32 used classification accuracy to measure, while 30 of them calculated the F1-score for their classification models. Precision and recall are also widely used, with both 29 papers reporting them. Moreover, the confusion matrix, ROC-AUC, task success rate are also applied to evaluate the model performance, although they were less common.

\subsection{Evaluation Metrics for Rule-Based Chatbots}

Since around 2018--2020, research shifted heavily towards retrieval-based chatbots, generative chatbots, and LLM-based agents, there are quite few papers talking about the evaluation metrics for pure rule-based chatbots. In the eight papers we collected, rule-based chatbots are generally evaluated from four aspects: Coverage, Robustness, User satisfaction.

\begin{description}
    \item[\textbf{Coverage}] Coverage is used to measure how much of the users' question space this chatbot can handle. An ideal chatbot should be able to response to most of the users' questions. And that's the reason why coverage should be considered as one of the core measurements for a chatbot's capability. To measure this, D'Ávila \cite{davila2018} introduced the KINO framework, which evaluates rule-based chatbot performance using two monitoring metrics---assertiveness (the ability to provide a valid response) and ambiguity (the degree of uncertainty caused by overlapping rules).

    \item[\textbf{Robustness}] Sometimes, users might click Enter accidentally without finishing typing their questions, or they might mistaken the spelling of some words. Therefore, robustness should also be considered as one of the important standards to evaluate a chatbot. D'Ávila \cite{davila2018} used out-of-scope (OOS) interactions in his KINO framework to refer to user inputs that cannot be matched to any existing conversational rule, intent, or response pattern. A high proportion of OOS inputs suggests reduced conversational robustness, whereas a lower OOS rate indicates that the chatbot can successfully handle a broader range of user requests within its intended domain.

    \item[\textbf{User satisfaction}] As a chatbot designed to provide users with convenience, user satisfaction is absolutely an important standard to evaluate a chatbot. To measure this, Papakostas et al.\ \cite{papakostas2024} used a questionnaire administered to 50 undergraduate students. The questionnaire covered general user satisfaction, query accuracy, response time, and personalized guidance. By directly questioning users, the authors can get a grasp of the actual user satisfaction conveniently. However the problem is, surveying a sufficiently large number of people might be a costly and labor-intensive work, whereas surveying just a small number of respondents may lead to biased conclusion. Moreover, Dhinagaran et al.\ \cite{dhinagaran2022} provided another way to evaluate user satisfaction. They use retention (the proportion of users who continue using the chatbot over time) and number of interactions (the number of rounds a user need before their query is resolved) to measure the user satisfaction. If a user finds the chatbot helpful, they will keep using it, and if the chatbot can solve the user's problem efficiently, the user will get satisfied. Therefore, these two are indeed reasonable evaluation metrics.
\end{description}

\subsection{Evaluation Metrics for Retrieval-Based Chatbots}

For retrieval-based chatbots, Recall@k (measures whether the correct answer appears within the top $k$ retrieved responses), MRR (measures how highly the first correct answer is ranked), and Precision@k (measures how many of the top $k$ retrieved results are actually relevant) are most commonly used. Metrics like ROUGE and BERTScore are also frequently mentioned, but they are more commonly used for retrieval-augmented generation (RAG) chatbots rather than pure retrieval-based chatbots, so they are not the focus here.

Interestingly, another paper we found pointed out the limitation of the previously mentioned metrics. Liu et al.\ \cite{liu2021meta} performed a meta-evaluation. Instead of evaluating chatbots, they evaluate the classic evaluation metrics themselves. Their finding indicates that most commonly used metrics do not strongly predict what users actually prefer. A model with high Recall@k, Precision@k, etc.\ can still provide poor user experience. Inspired by their findings, we propose that a complete evaluation for retrieval based chatbots should consist of two complementary literature streams: Response Selection Evaluation (classic evaluation) and Conversational Evaluation Metrics (user satisfaction). In this way the developer can have a comprehensive grasp of the model performance.

\section{Conclusion}
This literature review examined existing university chatbot applications, 
the technical approaches used in chatbot development, 
and the evaluation methods adopted in previous studies.

The review found that university chatbots are widely applied in areas 
such as student enquiries, academic support, admissions, and administrative services. 
Different types of chatbots can support freshmen in different aspects.
Among the different approaches, each has its own strengths and limitations.
Hybrid architecture is a useful way to leverage the strengths and overcome the shortcomings of different modules, and it has been widely applied by many researchers.
The review also analyzed the commonly used evaluation metrics for different chatbot components,
and concluded that a comprehensive assessment should include evaluation from different aspects.

Overall, this literature review analyzes chatbots developed by previous researchers 
and provides a useful starting point for building FAQ chatbots.


\clearpage
\begin{thebibliography}{99}

\bibitem{martinez2024}
S. Martinez-Requejo, E. J. García, S. R. Duarte, J. R. Lázaro, E. P. Sanz, and G. M. Vivas, ``AI-driven student assistance: Chatbots redefining university support,'' in \textit{INTED2024 Proceedings}, pp. 617--625, IATED, 2024.

\bibitem{sharma2019}
R. M. Sharma, ``Chatbot based college information system,'' \textit{RESEARCH REVIEW International Journal of Multidisciplinary}, vol. 4, no. 3, pp. 109--112, 2019.

\bibitem{susanna2020}
M. C. L. Susanna, R. Pratyusha, P. Swathi, P. R. Krishna, and V. S. Pradeep, ``College enquiry chatbot,'' \textit{International Research Journal of Engineering and Technology (IRJET)}, vol. 7, no. 3, pp. 784--788, 2020.

\bibitem{salve}
P. Salve, V. Patil, V. Gaikwad, and G. Wadhwa, ``College enquiry chat bot,'' \textit{International Journal on Recent and Innovation Trends in Computing and Communication}, vol. 5, no. 3, n.d.

\bibitem{dhandayuthapani2022}
V. B. Dhandayuthapani, ``A proposed cognitive framework model for a student support chatbot in a higher education institution,'' \textit{International Journal of Advanced Networking and Applications}, vol. 14, no. 2, pp. 5390--5395, 2022.

\bibitem{nourbakhsh2017}
A. Nourbakhsh and S. Yaser, ``Chatbots in education: A review of the state-of-the-art and future directions,'' \textit{Educational Technology Research and Development}, vol. 65, no. 5, pp. 1079--1104, 2017.

\bibitem{qian2019}
Y. Qian, \textit{How Do Freshmen Find Information Orienting Themselves in Their First Year of College? An Investigation of the Information Seeking Behaviors of First-Year Undergraduate Students}, 2019.

\bibitem{jones2008}
S. Jones, \textit{Internet Goes to College: How Students Are Living in the Future with Today's Technology}. Diane Publishing, 2008.

\bibitem{raval2020}
H. Raval, ``Limitations of existing chatbot with analytical survey to enhance the functionality using emerging technology,'' \textit{International Journal of Research and Analytical Reviews (IJRAR)}, vol. 7, no. 2, 2020.

\bibitem{khin2020}
N. N. Khin and K. M. Soe, ``University chatbot using artificial intelligence markup language,'' in \textit{Proc. IEEE Conf. Computer Applications (ICCA)}, pp. 1--5, Feb. 2020.

\bibitem{nguyen2021}
T. T. Nguyen, A. D. Le, H. T. Hoang, and T. Nguyen, ``NEU-chatbot: Chatbot for admission of National Economics University,'' \textit{Computers and Education: Artificial Intelligence}, vol. 2, Art. no. 100036, Jan. 2021.

\bibitem{alghadhban2020}
D. Al-Ghadhban and N. Al-Twairesh, ``Nabiha: An Arabic dialect chatbot,'' \textit{International Journal of Advanced Computer Science and Applications}, vol. 11, no. 3, pp. 452--459, 2020.

\bibitem{dinesh2021}
T. Dinesh, M. R. Anala, T. T. Newton, and G. R. Smitha, ``AI bot for academic schedules using Rasa,'' in \textit{Proc. International Conference on Innovative Computing, Intelligent Communication and Smart Electrical Systems (ICSES)}, pp. 1--6, Sep. 2021.

\bibitem{lopez2022}
T. Lopez and M. Qamber, \textit{The Benefits and Drawbacks of Implementing Chatbots in Higher Education: A Case Study for International Students at Jönköping University}, M.S. thesis, Jönköping International Business School, Jönköping University, Sweden, 2022.

\bibitem{ashraf2024}
N. Ashraf, R. S. Ahmad, S. Bano, H. M. Azeem, and S. Naz, ``Enhancing MBTI personality prediction from text data with advance word embedding technique,'' \textit{VFAST Transactions on Software Engineering}, vol. 12, no. 3, pp. 35--43, 2024.

\bibitem{yang2026}
X. Yang, ``Generative language model for chatbots based on TF-IDF personality feature recognition and classification,'' \textit{Modelling and Simulation in Engineering}, vol. 2026, Art. ID 2635948, Mar. 2026, doi: \href{https://doi.org/10.1155/mse/2635948}{10.1155/mse/2635948}.

\bibitem{mccallum1998}
A. McCallum and K. Nigam, ``A comparison of event models for Naive Bayes text classification,'' in \textit{Proc. AAAI Workshop on Learning for Text Categorization}, pp. 41--48, 1998.

\bibitem{setyaningsih2021}
E. R. Setyaningsih and I. Listiowarni, ``Categorization of exam questions based on Bloom taxonomy using Naive Bayes and Laplace smoothing,'' in \textit{Proc. 3rd East Indonesia Conference on Computer and Information Technology (EIConCIT)}, Surabaya, Indonesia, pp. 330--333, 2021, doi: \href{https://doi.org/10.1109/EIConCIT50028.2021.9431862}{10.1109/EIConCIT50028.2021.9431862}.

\bibitem{acito2024}
F. Acito, ``Logistic regression,'' in \textit{Predictive Analytics with KNIME}, Cham, Switzerland: Springer Nature, pp. 125--167, 2024, doi: \href{https://doi.org/10.1007/978-3-031-45630-5_7}{10.1007/978-3-031-45630-5\_7}.

\bibitem{li2009}
G. Li, S. Xing, and H. Xue, ``Comparison on pattern analysis performance of SVM and RVM based on RBF kernel,'' \textit{Application Research of Computers}, vol. 26, no. 5, pp. 1782--1784, 2009 (in Chinese).

\bibitem{liu2010}
Z. Liu, X. Lv, K. Liu, and S. Shi, ``Study on SVM compared with the other text classification methods,'' in \textit{2010 Second International Workshop on Education Technology and Computer Science}, Wuhan, China, pp. 219--222, 2010, doi: \href{https://doi.org/10.1109/ETCS.2010.248}{10.1109/ETCS.2010.248}.

\bibitem{mikael2025}
K. Mikael, C. Öz, T. A. Rashid, and G. S. Nariman, ``A hybrid chatbot model for enhancing administrative support in education: Comparative analysis, integration, and optimization,'' \textit{IEEE Access}, vol. 13, pp. 50741--50754, 2025, doi: \href{https://doi.org/10.1109/ACCESS.2025.3552501}{10.1109/ACCESS.2025.3552501}.

\bibitem{salton1988}
G. Salton and C. Buckley, ``Term-weighting approaches in automatic text retrieval,'' \textit{Information Processing \& Management}, vol. 24, no. 5, pp. 513--523, 1988, doi: \href{https://doi.org/10.1016/0306-4573(88)90021-0}{10.1016/0306-4573(88)90021-0}.

\bibitem{robertson2009}
S. Robertson and H. Zaragoza, ``The probabilistic relevance framework: BM25 and beyond,'' \textit{Foundations and Trends in Information Retrieval}, vol. 3, no. 4, pp. 333--389, 2009, doi: \href{https://doi.org/10.1561/1500000019}{10.1561/1500000019}.

\bibitem{sakata2019}
W. Sakata, T. Shibata, R. Tanaka, and S. Kurohashi, ``FAQ retrieval using query-question similarity and BERT-based query-answer relevance,'' \textit{arXiv preprint arXiv:1905.02851}, 2019.

\bibitem{reimers2019}
N. Reimers and I. Gurevych, ``Sentence-BERT: Sentence embeddings using Siamese BERT-networks,'' in \textit{Proc. 2019 Conference on Empirical Methods in Natural Language Processing and the 9th International Joint Conference on Natural Language Processing (EMNLP-IJCNLP)}, Hong Kong, China, pp. 3982--3992, 2019, doi: \href{https://doi.org/10.18653/v1/D19-1410}{10.18653/v1/D19-1410}.

\bibitem{karpukhin2020}
V. Karpukhin, B. Oguz, S. Min, P. Lewis, L. Wu, S. Edunov, D. Chen, and W.-T. Yih, ``Dense passage retrieval for open-domain question answering,'' in \textit{Proc. 2020 Conference on Empirical Methods in Natural Language Processing (EMNLP)}, pp. 6769--6781, 2020, doi: \href{https://doi.org/10.18653/v1/2020.emnlp-main.550}{10.18653/v1/2020.emnlp-main.550}.

\bibitem{jose2025}
T. A. Jose, Abhishek, A. P. Patil, A. G. S, and B. Bharath, ``Hybrid chatbot for educational institutions,'' \textit{International Journal of Recent Trends in Multidisciplinary Research}, vol. 5, no. 6, pp. 151--161, Nov.--Dec. 2025, doi: \href{https://doi.org/10.59256/ijrtmr.20250506020}{10.59256/ijrtmr.20250506020}.

\bibitem{davila2018}
T. C. D'Ávila, \textit{KINO: An Approach for Rule-Based Chatbot Development, Monitoring and Evaluation}, Master's thesis, Federal University of Minas Gerais, Brazil, 2018.

\bibitem{papakostas2024}
C. Papakostas, C. Troussas, A. Krouska, and C. Sgouropoulou, ``A rule-based chatbot offering personalized guidance in computer programming education,'' in \textit{Intelligent Tutoring Systems}, Springer, 2024, doi: \href{https://doi.org/10.1007/978-3-031-63031-6_22}{10.1007/978-3-031-63031-6\_22}.

\bibitem{dhinagaran2022}
D. A. Dhinagaran, L. Martinengo, M. H. R. Ho, S. Joty, T. Kowatsch, R. Atun, and L. Tudor Car, ``Designing, developing, evaluating, and implementing a smartphone-delivered, rule-based conversational agent (DISCOVER): Development of a conceptual framework,'' \textit{JMIR mHealth and uHealth}, vol. 10, no. 10, p. e38740, 2022, doi: \href{https://doi.org/10.2196/38740}{10.2196/38740}.

\bibitem{liu2021meta}
Z. Liu, K. Zhou, and M. L. Wilson, ``Meta-evaluation of conversational search evaluation metrics,'' \textit{ACM Transactions on Information Systems}, vol. 39, no. 4, Art. 45, 2021, doi: \href{https://doi.org/10.1145/3445029}{10.1145/3445029}.

\end{thebibliography}

\end{document}
